\chapter*{KẾT LUẬN}
\addcontentsline{toc}{chapter}{Kết luận}

Phần việc TV2 đã hoàn thành hai cải tiến cho thuật toán NSGA-II giải bài toán MOGVRPTW của nhóm. Khởi tạo
thông minh tạo quần thể ban đầu từ thứ tự láng giềng gần nhất, thứ tự theo khung giờ và hoán vị ngẫu
nhiên theo tỉ lệ cấu hình, có biến thể và khử trùng. Tìm kiếm cục bộ cải thiện từng con bằng 2-opt và
Or-opt trong từng tuyến, chấm mọi ứng viên bằng đúng bộ đánh giá của thuật toán, có hai tiêu chí chấp
nhận, hai chính sách xử lý lệch biên tuyến và ngân sách đánh giá cố định. Cả hai đã được gắn vào
\code{nsga2.py} của nhóm, chạy được cùng các cải tiến của TV3 và TV4, và khi tắt cờ thì kết quả trùng từng
bit với thuật toán gốc. Toàn bộ 432 test của dự án đều pass.

Thực nghiệm trên bốn instance với 10 seed cho thấy khởi tạo thông minh tăng HV có ý nghĩa thống kê trên cả
bốn instance mà không tốn thêm lần đánh giá nào. Tìm kiếm cục bộ đứng riêng ít hiệu quả vì phần lớn bước
đi bị lệch biên tuyến với cách cắt \code{capacity\_time}; chính sách \code{evaluate\_actual} khắc phục được
một phần. Kết hợp hai cải tiến cho HV cao nhất ở cùng 200 thế hệ trên cả bốn instance, và vẫn tốt hơn thuật
toán gốc được cấp cùng số lần đánh giá trên C102-25 và RC201-100.

Hướng tiếp theo là đánh giá tăng dần để giảm chi phí tìm kiếm cục bộ, thêm bước đi giữa các tuyến, chọn
lại chính sách lệch biên mặc định và mở rộng thực nghiệm ra đủ chín instance của bài báo.
